\scanpage{63}
% archive foliation: 62
On \deleted{\unsure{veut}} se demande s'il est représentable, i.e. on demande la valeur en $\unsure{T}$ du \uline{foncteur adjoint} (au sens de Kan) du foncteur $\underline W$ des $S$-préschémas en $\unsure{T}$-préschémas. Une solution du Pb concret posé par le foncteur en $T$ ci-dessus est par définition un couple
\[
X\quad\text{($S$-préschéma.)}\qquad \underline W(X)\xrightarrow{\alpha}\mathfrak X
\]
tel que pour tout $T$, l'application naturelle
\[
\mathrm{Hom}_S(T,X)\xrightarrow{\sim}\mathrm{Hom}_{\underline W(S)}(\underline W(T),\mathfrak X)
\]
\deleted{soit un} définie par \unsure{factorisation}
\diagram{\begin{tikzcd}
\underline W(T)\arrow[r,"\underline W(u)"]&\underline W(X)\arrow[r,"\alpha"]&\mathfrak X
\end{tikzcd}}
soit \uline{bijective}. Comme le foncteur $(\ast)$ en $T$ est de ``nature locale'', il suffit d'ailleurs de vérifier la bijectivité pour $T$ affine \deleted{\unsure{à}} \inserted{sur un ouvert} affine de $S$, p. ex.

Désignons par $\underline W^a$ le foncteur adjoint de $\underline W$ (défini sur une ss-catégorie de la catégorie des $\unsure{T}$-préschémas)
\margin{$\mathrm{Hom}_{S^a}(T^a,\mathfrak X)\simeq\mathrm{Hom}_S(T,\mathfrak X$}
